%! TEX root = ./main.tex


%%%%%%%%%%%%%%%%%%%%%%%%%%%%%%%%%%%%%%%%%%%%%%%%%%%%%%%%%%%%%%%%%%%%%%%%%%%%%%%%%%%%%%%%%%%%%%%%%%%%
%%%%%%%%%%%%%%%%%%%%%%%%%%%%%%% SECTION 3 %%%%%%%%%%%%%%%%%%%%%%%%%%%%%%%%%%%%%%%%%%%%%%%%%%%%%%%%%%
%%%%%%%%%%%%%%%%%%%%%%%%%%%%%%%%%%%%%%%%%%%%%%%%%%%%%%%%%%%%%%%%%%%%%%%%%%%%%%%%%%%%%%%%%%%%%%%%%%%%
%%%%%%%%%%%%%%%%%%%%%%%%%%%%%%%%%%%%%%%%%%%%%%%%%%%%%%%%%%%%%%%%%%%%%%%%%%%%%%%%%%%%%%%%%%%%%%%%%%%%

\section{\texorpdfstring{Proof of the main result: case $\beta \in \Q$}{}}

\label{sec:new beta map}

We can use the dynamical system realization of the chunk-by-chunk greedy map described in (\ref{eq:greedy map chunks as dynamical system}) to generate chunks of a specific $\beta$-expansion of a given $s \in [0,1]$ by reading increasingly long prefixes of its greedy binary expansion. The generated $\beta$-expansion will then be proven to satisfy (\ref{eq:inequality to be proven}). The convertion prefix-to-chunk appears to be an important tool to control the behavior of the dynamical system, so that it indeed generates a $\beta$-expansion. Let $s \in [0,1]$ and $N \in \N$. When presented with increasingly long prefixes of $\gf_2(s)$ of respective lengths $\Sigma_i \in \N$, the dynamical system converts $\gf_2(s)_{1: \Sigma_i}$ to $\xf_{1: Ni}$, for $i \in \N$, where $\xf \in \Sigma_\beta(s)$. We will choose the $\Sigma_i$ to be such that $\limi{i} \frac{\Sigma_i}{Ni} = \log_2(\beta)$, which leads to (\ref{eq:inequality to be proven}) for $\beta \in \Q$. Specifically, we shall see that the dynamical system is realized by an effective algorithm. Figure \ref{fig:principle conversion greedy binary to beta exp} illustrates this process.
\begin{figure}[H]
	\centering
	\tikzmath{\x=3;}
	\tikzmath{\y1=0.25;}
	\tikzmath{\y2=0.5;}
	\tikzmath{\S=0.93;}
	\begin{tikzpicture}[scale = 0.8]
		\node[anchor = west] at (0,0.1){{\small $\gf_2(s) = 01001|0001|111|1001|000|1101|010|0000|001|1000|011|\ldots$}};
		\node[anchor = west] at (0,-\x){{\small$\xf = \underset{N}{000100}|\underset{N}{101000}|\underset{N}{001000}|101000|100000|101000|100000|100000|100000|010001|000010|\ldots$}};
		\draw[thick, ->] (2,-1.7) -- (2.9*\S-1, -\x + \y2);
		\draw[thick, ->] (4*\S,-1) -- (4.6*\S-0.7, -\x + \y2);
		\draw[thick, ->] (5*\S,-0.4) -- (6.1*\S-0.8, -\x + \y2);
		\draw [decorate,decoration={brace,amplitude=5pt,mirror}]
  (1.9,-1.5) -- (3.1,-1.5) node[midway, yshift = -0.9em]{\scriptsize$\Sigma_1$};
  \draw [decorate,decoration={brace,amplitude=5pt,mirror}]
  (1.9,-0.8) -- (4.2,-0.8) node[midway, yshift = -0.9em]{\scriptsize$\Sigma_2$};
  \draw [decorate,decoration={brace,amplitude=5pt,mirror}]
  (1.9,-0.22) -- (5,-0.22) node[midway, yshift = -0.9em]{\scriptsize$\Sigma_3$};
  	\path (5.5,-1.2) -- node[auto=false]{\ldots} (6.5,-1.2);
	\end{tikzpicture}
	\caption{Conversion of the greedy binary expansion $\gf_2(s)$ to a $\beta$-expansion $\xf$ of similar algorithmic complexity.}
	\label{fig:principle conversion greedy binary to beta exp}
\end{figure}

The chunks of the $\beta$-expansion $\xf$ are generated by the following iterative process. Suppose that at step $i \in \N_0$, $\xf_{1:Ni}$ has been generated from $\gf_2(s)_{1:\Sigma_i}$. Now, we read a next chunk of $\gf_2(s)$, of size $\Sigma_{i+1}-\Sigma_i$, so that we have at our disposal $\xf_{1:Ni}$ and $\gf_2(s)_{1:\Sigma_{i+1}}$, and we generate a chunk of $\xf$ of size $N$ by choosing the lexicographically maximal binary sequence $x \in \{0,1\}^N$ that satisfies 
\begin{equation}\label{eq:reducing the error between x beta and greedy binary}
	\delta_2(\gf_2(s)_{1:\Sigma_{i+1}}) - \frac{\beta^{-N(i+1)}}{\beta - 1} \leq \delta_\beta(\xf_{1:Ni}x) \leq \delta_2(\gf_2(s)_{1:\Sigma_{i+1}}),
\end{equation}
and take $\xf_{1:N(i+1)} := \xf_{1:Ni} x$. Using Lemma \ref{lem:factorisation of beta expansion through concatenation two}, one can rewrite (\ref{eq:reducing the error between x beta and greedy binary}) as
\begin{equation}\label{eq:first step towards x greedy exp of R i}
	\delta_2(\gf_2(s)_{1:\Sigma_{i+1}}) - \frac{\beta^{-N(i+1)}}{\beta - 1} \leq \delta_\beta(\xf_{1:Ni}) + \beta^{-Ni} \delta_\beta(x) \leq \delta_2(\gf_2(s)_{1:\Sigma_{i+1}}).
\end{equation}
Using again Lemma \ref{lem:factorisation of beta expansion through concatenation two}, we have
\begin{equation}\label{eq:factorisation binary expansion}
	\delta_2(\gf_2(s)_{1:\Sigma_{i+1}}) = \delta_2\left(\gf_2(s)_{1:\Sigma_{i}}\right) + 2^{-\Sigma_i}\delta_2\left(\gf_2(s)_{\Sigma_i+1:\Sigma_{i+1}}\right).
\end{equation}
By defining $R_0 = 0$, $\Sigma_0 = 0$,
\begin{equation}\label{eq:R i definition}
	R_i := \beta^{Ni}\left(\delta_2(\gf_2(s)_{1:\Sigma_{i}}) - \delta_\beta(\xf_{1:Ni})\right), \ \text{for} \ i \in \N,
	\tag{\text{$\mathcal{R}$}}
\end{equation}
and
\begin{equation}\label{eq:s i to generate beta expansion}
	s_i := \frac{\beta^{Ni}}{2^{\Sigma_i}}\delta_2(\gf_2(s)_{\Sigma_i+1:\Sigma_{i+1}}), \ \text{for} \ i \in \N_0,
	\tag{\text{$\mathcal{S}$}}
\end{equation}
(\ref{eq:first step towards x greedy exp of R i}) is equivalent to 
\begin{equation}\label{eq:greedy exp for x}
	R_i + s_i - \frac{\beta^{-N}}{\beta-1} \leq \delta_\beta(x) \leq R_i + s_i.
\end{equation}
$x \in \{0,1\}^N$ is then the lexicographically maximal sequence satisfying (\ref{eq:greedy exp for x}), which exactly corresponds to the first $N$ bits of the greedy $\beta$-expansion of $R_i+s_i$, i.e.
\begin{equation}
	x := d_\beta^N(R_i + s_i).
\end{equation}
Therefore, $\xf$ is iteratively defined as 
\begin{equation}
	\xf = \coprod_{i=0}^\infty \grpref\left(R_i + s_i\right).
\end{equation}
This exatcly matches (\ref{eq:inferring chunks of greedy expansion from R i E beta}), and
hence, the successive values taken by the variable $R_i$ correspond to the evolution of the dynamical system (\ref{eq:greedy map chunks as dynamical system}) given by
\begin{equation}
	\tag{\text{$E^N_\beta$}}
	R_{i+1} = G_\beta^N(R_i + s_i),
\end{equation}
when presented with input $\ssf =s_0s_1\ldots$ defined by (\ref{eq:s i to generate beta expansion}). 
However, one must pay attention to the fact that if $R_i + s_i \notin I_\beta$, then (\ref{eq:greedy map chunks as dynamical system}) is not well-defined as $G_\beta$ is defined on $I_\beta$ only. Using the interpretation in terms of a dynamical system, we establish conditions on $N,\Sigma_i$ so that, given $\beta$, (\ref{eq:greedy map chunks as dynamical system}) is well-defined, then we prove that the above procedure produces a valid $\beta$-expansion, and finally we study the algorithmic complexity of the delivered $\beta$-expansion.







% \begin{figure}[H]
% 	\centering
% 	\includegraphics[width=0.9\textwidth]{images/conversion_greedy_binary_to_beta_low_complex.png}
% 	% \caption{Principle of conversion from the greedy binary expansion $g_2(s)$ to a $\beta$-expansion $\xf_\beta(s)$ with similar algorithmic complexity.}
% 	% \label{fig:principle conversion greedy binary to beta exp}
% \end{figure}

% Before formally introducing the dynamical system, we show that the greedy map $G_\beta$ can be interpreted as a bit-shifting and a bit-extraction device. This vantage point is helpful in understanding the operational aspects of the dynamical system. Let $s \in [0,1]$ and denote its greedy expansion by $\xf$. Then, one has 
% \begin{equation}
% 	G_\beta(s) = G_\beta\left(\sum_{i=1}^\infty \xf_i \beta^{-i}\right) = \beta \left(\sum_{i=1}^\infty \xf_i \beta^{-i}\right) - \xf_1 = \sum_{i=1}^\infty \xf_{i+1} \beta^{-i}.
% \end{equation}
% which shows that application of $G_\beta$ simply shifts the sequence $x$ by one position to the left. It now follows directly that iterating $G_\beta$ $n$ times effects an left-shift by $n$ positions according to 
% \begin{equation}\label{eq:n iterations of greedy map is same as shifting n bits}
% 	G^n_\beta(s) = \sum_{i=1}^\infty \xf_{i+n} \beta^{-i}.
% \end{equation}
% Moreover, we can also interpret the greedy map as a bit-extraction device, namely by noting that
% \begin{equation}
% 	s - \frac{1}{\beta} G_\beta(s) = \sum_{i=1}^\infty \xf_i \beta^{-i} - \frac{1}{\beta}\sum_{i=1}^\infty \xf_{i+1} \beta^{-i} = \frac{\xf_1}{\beta}.
% \end{equation}

% Likewise, the first $n$ bits of the greedy $\beta$-expansion of $s$ can be extracted as
% \begin{equation}\label{eq:n iterations of greedy map is same as reading the first n bits}
% 	s - \frac{1}{\beta^n} G^n_\beta(s) = \sum_{i=1}^\infty \xf_i \beta^{-i} - \frac{1}{\beta^n}\sum_{i=1}^\infty \xf_{i+n} \beta^{-i} = \sum_{i=1}^n\xf_i \beta^{-i} = \delta_\beta(\xf_{1: n}).
% \end{equation}



% We will later prove that $\dom \beta$ is not empty for every $\beta \in (1,2)$.
% Accordingly, we define
% \begin{equation}\label{eq:definition of the low complex beta expansion}
% 	\xf_\beta(s)_{1:Ni} := \DD_\Theta(s_{0:i-1}), \ \ \text{for all} \ i \in \N_0,
% \end{equation}
% with $s_i$ satisfying (\ref{eq:s i to generate beta expansion}).

% We turn to the description of the ideas underlying the conception of the dynamical system just introduced. We suppose that $S_0 = s \in [0,1]$. 
% \begin{enumerate}
% 	\item We first give an interpretation of the sequence $S_i$. By iteration of (\ref{eq:definition of K beta}), one gets $S_i = G_2^{\sum_{j=0}^{i-1} M_j}(S_0)= G_2^{\Sigma_{i}}(s)$, for all $i \in \N$. As $G_2$ performs bitshifting in the greedy binary expansion, $S_i$ encodes the greedy binary expansion of $s$, with the first $\Sigma_{i} = \sum_{j=0}^{i-1} M_j$ bits removed. In Figure \ref{fig:principle conversion greedy binary to beta exp}, $S_i$ acts as a pointer indicating the starting locations of the chunks of size $M_{i}$ of the greedy binary expansion of $s$.
% 	\item We now give an interpretation of the sequence $S'_i$. Recall that $S'_i = S_{i} - \frac{1}{2^{M_{i}}}G_2^{M_{i}}(S_{i})$. As $s \mapsto s - \frac{1}{2^n}G^n_2(s)$ extracts the first $n$ bits of the greedy binary expansion of $s$, we deduce that $S'_i = S_{i} - \frac{1}{2^{M_{i}}}G_2^{M_{i}}(S_{i})$ encodes the first $M_i$ bits of the greedy binary expansion of $S_i$. Hence, $S'_i$ encodes the chunks of greedy binary expansion of $s$, with respective lengths $M_i$, and respective starting locations $\Sigma_i$.
% 	\item We move to the interpretation interpretation of the sequences $R_i$ and $R'_i$. $R_i$ is designed so that
% 	\begin{equation}
% 		R_i = \beta^{Ni}\left(\delta_2(\gf_2(s)_{1:\Sigma_i}) - \delta_\beta(X_i)\right),
% 	\end{equation}
% 	and can be interpreted as the approximation error between the $Ni$-prefix of $\xf_\beta(s)$ and the approximation of $s$ by the $\Sigma_i$-truncation of its greedy binary expansion. Further, $R'_i$ is designed so that
% 	\begin{equation}
% 		R'_i = \beta^{Ni}\left(\delta_2(\gf_2(s)_{1:\Sigma_{i+1}}) - \delta_\beta(X_i)\right),
% 	\end{equation}
% 	and can be interpreted as the approximation error between the $Ni$-prefix of $\xf_\beta(s)$ and the approximation of $s$ by the $\Sigma_{i+1}$-truncation of its greedy binary expansion. The goal of the dynamical system is to reduce this approximation error by adding a chunk of size $N$ to the $Ni$-prefix of $\xf_\beta(s)$.
	
% 	The goal of the dynamical system is to constantly aim at $\delta_2(\gf_2(s)_{1:\Sigma_i})$ by adding bits to $X_i$, which is corresponds exactly to a controlled dynamical system in control theory. On Figure A is represented the block diagram of this controlled system.
% \end{enumerate}

% We now show that, indeed, (\ref{eq:definition of K beta}) is well-defined under suitable choices of $N$ and $\Sigma_i$. 
% Note that (\ref{eq:definition of K beta}) is well-defined is equivalent to say that (\ref{eq:reducing the error between x beta and greedy binary}) is satisfied, which, by the above discussion, implies that $\DD_\Theta(\ssf)$ is a $\beta$-expansion of $s$. We will prove this statement formally in the next subsection.

\subsection{The dynamical system is well-defined}



For (\ref{eq:greedy map chunks as dynamical system}) to be well-defined, we must have $R_i + s_i \in I_\beta$ for all $i \in \N_0$. We define
\begin{equation}
	\dom E_\beta^N := \left\{ \ssf \in I_\beta^\omega : R_i + s_i \in I_\beta, \ \forall i \in \N_0  \right\}.
\end{equation}
We will use the following result on the greedy map $G_\beta$, inspired from a note\footnote{On page 317: ``Notice that for each $x \in [0, \lfloor \beta\rfloor/(\beta-1)]$ there exists an integer $n_0 = n_0(x)$ such that for all $n \geq n_0$ one has $T_\beta^n(x) \in [0,1)$.''} in \cite[Section 1]{dajani2002greedy}, which states for all $t \in I_\beta$ and large enough $N \in \N$, $G_\beta^N(t) \in [0,1]$. This result will allow to tune $N$ so that $0 \leq R_i \leq 1$, and we will then tune the $\Sigma_i$ so that $0 \leq s_i \leq \frac{\beta}{2(\beta - 1)} - 1$. Altogether, this indeed yields $0 \leq R_i + s_i \leq \frac{\beta}{2} \frac{1}{\beta-1} < \frac{1}{\beta-1}$. We start with the result on the greedy map.
\begin{lemma}\label{ineq:greedy beta expansion}
	Let $\beta \in (1,2)$, $r \in I_\beta$, and set $N_\beta(r) =
		-\log_\beta\left(\frac{1-(\beta-1)r}{2 - \beta}\right)$. Then, for $n \in \N$, it holds that $n > N_\beta(r)$ if and only if
	\begin{equation}
		G^n_\beta(r) \in [0,1).
	\end{equation}
\end{lemma}
\begin{proof} Fix $\beta \in (1,2)$. The proof is divided two regimes for $r \in I_\beta$. First, note that if $r \in [0,1)$, then
	\begin{equation}
		G_\beta(r) \overset{(\ref{eq:greedy map})}{=} \begin{cases}
		\beta r, \ \text{if} \ r < \frac{1}{\beta},\\
		\beta r - 1, \ \text{if} \ r \geq \frac{1}{\beta}, 
	\end{cases} \geq \begin{cases}
		\beta \times 0, \ \text{if} \ r < \frac{1}{\beta},\\
		\beta \times \frac{1}{\beta} -1, \ \text{if} \ r \geq \frac{1}{\beta}, 
	\end{cases} = 0.
	\end{equation}
	and
	\begin{equation}
		G_\beta(r) \overset{(\ref{eq:greedy map})}{=} \begin{cases}
		\beta r, \ \text{if} \ r < \frac{1}{\beta},\\
		\beta r - 1, \ \text{if} \ r \geq \frac{1}{\beta}, 
	\end{cases} < \begin{cases}
		\beta \times \frac{1}{\beta}, \ \text{if} \ r < \frac{1}{\beta},\\
		\beta \times 1 -1, \ \text{if} \ r \geq \frac{1}{\beta}, 
	\end{cases} \leq 1.
	\end{equation}
	Therefore, $G_\beta([0,1)) \subseteq [0,1)$. Now, fix $x \in I_\beta \backslash [0,1) = \left[1,\frac{1}{\beta-1}\right]$, and take $m \in \N$ to satisfy $G^m_\beta(r) \geq 1$. As $G_\beta([0,1)) \subseteq [0,1)$, we must have for all $i \leq m$, that $G^i_\beta(r) \geq 1 \geq \frac{1}{\beta}$. Then, one gets
	\begin{equation}\label{eq:value of m to satisfy non greedy condition}
		1 \leq G^m_\beta(r) \overset{(\ref{eq:greedy map})}{=} (t \mapsto \beta t - 1)^m(r) = \beta^m r - \sum_{i=0}^{m-1} \beta^i = \beta^m r - \frac{\beta^m-1}{\beta-1} \cdot
	\end{equation}
	This is satisfied if and only if
	\begin{equation}
		m \leq -\log_\beta\left(\frac{1-(\beta-1)r}{2 - \beta}\right).
	\end{equation}
\end{proof}
\noindent We now move to the statement allowing to control $R_i$.
\begin{lemma}\label{eq:J beta is in dom E beta N}
	Let $\beta \in (1,2)$. If $N \geq \log_\beta(2)$ then $J_\beta^\omega \subset \dom E_\beta^N$, with 
	\begin{equation}\label{eq:definition J beta}
		J_\beta := \left[0, \frac{\beta}{2(\beta - 1)} - 1 \right] \neq \varnothing.
	\end{equation}
\end{lemma}
\begin{proof}
	Fix $\beta \in (1,2)$. We start by showing that $J_\beta$ is nonempty.
	The function $f : (1,2] \to \R$ defined as $f(\beta) = \frac{\beta/2}{\beta-1}$ for all $\beta \in (1,2]$ satisifies $f'(\beta) = -\frac{1}{2(\beta-1)^2} < 0$ for all $\beta \in (1,2]$ and $f(2) = 1$. Hence, $f$ is decreasing on $(1,2]$ and $f(2) = 1$, so $ \frac{\beta}{2(\beta - 1)} - 1 = f(\beta) - 1 > 0$ for all $\beta \in (1,2)$, yielding $J_\beta \neq \varnothing$.

	Now, let $\ssf = s_0s_1s_2\ldots \in J_\beta^\omega$. We suppose that $N \geq \log_\beta(2)$, and prove by induction
	\begin{equation}\label{eq:condition Ri si}
		R_i + s_i \in I'_\beta := \left[0, \frac{\beta/2}{\beta-1}\right] \subset \left[0, \frac{1}{\beta-1}\right] = I_\beta, \ \text{for all}\ i \in \N_0.
	\end{equation}
	We start by showing that $R_0 + s_0 \in I'_\beta$. Indeed, from the definition of $E_\beta^N$, $R_0 = 0$. Moreover, as $s_0 \in J_\beta$, one has
	\begin{equation}
		R_0 + s_0 = s_0 \in J_\beta \subset I'_\beta.
	\end{equation}
	We continue by proving that if $R_{i-1} + s_{i-1} \in I'_\beta$, then $R_{i} + s_i \in I'_\beta$, for all $i \in \N$. More precisely, we show that if $R_{i-1} + s_{i-1} \in I'_\beta$, then $R_i \in [0,1]$, which combined with $s_i \in J_\beta$ yields $R_i + s_i \in I'_\beta$, for all $i \in \N$. Fix $i \in \N$ and suppose that $R_{i-1} + s_{i-1}\in I'_\beta$. It follows from (\ref{eq:greedy map chunks as dynamical system}) that
	\begin{equation}\label{eq:relation R i+1 R i s i}
		R_i = G_\beta^N\left(R_{i-1}+s_{i-1}\right).
	\end{equation}
	By Lemma \ref{ineq:greedy beta expansion}, 
	\begin{equation}
		G_\beta^n\left(R_{i-1} + s_{i-1}\right) \in [0,1]
	\end{equation}
	for all
	\begin{equation}\label{eq:condition on n for convergence of greedy expansion}
		n \geq - \log_\beta\left( \frac{1 - (\beta - 1)(R_{i-1} + s_{i-1})}{2- \beta} \right).
	\end{equation}
	From $R_{i-1} + s_{i-1} \leq \frac{\beta/2}{\beta-1}$, one calculates
	\begin{align}
		- \log_\beta\left( \frac{1 - (\beta - 1)(R_{i-1}+s_{i-1})}{2- \beta} \right) \leq - \log_\beta\left( \frac{1 - (\beta - 1)\frac{\beta/2}{\beta-1}}{2- \beta} \right)
		 = \log_\beta(2)
		 \leq N.
	\end{align}
	Then, $N$ satisfies (\ref{eq:condition on n for convergence of greedy expansion}), so
	\begin{equation}\label{eq:R i is in 0 1}
		R_{i} \overset{(\ref{eq:relation R i+1 R i s i})}{=} G^N_\beta(R_{i-1} + s_{i-1}) \osetlem{\ref{ineq:greedy beta expansion}}{\in} [0,1].
	\end{equation}
	By virtue of inductive reasonning, we conclude that $R_i + s_i \in I'_\beta$ for all $i \in \N_0$, so $\ssf \in \dom E_\beta^N$.
\end{proof}
\noindent We finally tune $\Sigma_i$ so that $\ssf \in J_\beta^\omega$.
\begin{lemma}\label{lem:s is in J beta}
	Let $\beta \in (1,2)$, $s \in [0,1]$, $N \in \N$, $(\Sigma_i)_{i \in \N_0}$, and $\ssf = s_0s_1\ldots$ be defined according to (\ref{eq:s i to generate beta expansion}). Then, if $\Sigma_0 = 0$ and
	\begin{equation}\label{eq:conditions on sigma i to work}
		\Sigma_i \geq Ni\log_2(\beta)+\log_2\left(\frac{\beta -1}{2 - \beta}\right) + 1, \ \text{for all} \ i \in \N,
	\end{equation}
	we have $\ssf \in J_\beta^\omega$.
\end{lemma}
\begin{proof}
	Let $\beta \in (1,2)$, $s \in [0,1]$, $N \in \N$, $(\Sigma_i)_{i \in \N_0}$ satisfying $\Sigma_0 = 0$ and (\ref{eq:conditions on sigma i to work}), and $\ssf = s_0s_1\ldots$ be defined by (\ref{eq:s i to generate beta expansion}). For $i \in \N_0$,
	\begin{equation}
		0 \leq s_i \overset{(\ref{eq:s i to generate beta expansion})}{=} \frac{\beta^{Ni}}{2^{\Sigma_i}} \delta_2(\gf_2(s)_{\Sigma_i+1:\Sigma_{i+1}}) \leq \frac{\beta^{Ni}}{2^{\Sigma_i}} \leq \frac{\beta/2}{\beta-1} - 1,
	\end{equation}
	where (a) follows from $\delta_2 : \{0,1\}^\omega \to [0,1]$, and (b) is by (\ref{eq:conditions on sigma i to work}).
\end{proof}
We wrap up the last three results in a Theorem that states precise values for $N$ and $(\Sigma_i)_{i \in \N_0}$ so that $\ssf \in \dom E_\beta^N$.
\begin{theorem}\label{thm:condition on the input sequences so that s is in dom E}
	Let $s \in [0,1]$, $\beta \in (1,2)$, $\Sigma_0 = 0$ and
	\begin{equation}\label{eq:final value parameters}\tag{\text{$\mathcal{P}$}}
		\begin{array}{rl}
			N &:= \lceil \log_\beta(2) \rceil,\\
			\Sigma_i &:= \lceil Ni\log_2(\beta)\rceil + \left\lceil \log_2\left(\frac{\beta -1}{2 - \beta}\right) + 1 \right\rceil, \ \text{for all} \ i \in \N.
		\end{array}
	\end{equation}
	Then, for all $\ssf = s_0s_1\ldots$ defined according to (\ref{eq:s i to generate beta expansion}), $\ssf \in \dom E_\beta^N$.
\end{theorem}
\begin{proof}
	$N = \lceil \log_\beta(2) \rceil \geq \log_\beta(2)$, so by Lemma \ref{eq:J beta is in dom E beta N}, $J_\beta^\omega \subset \dom E_\beta^N$. Moreover, $\Sigma_i := \lceil Ni\log_2(\beta)\rceil + \left\lceil \log_2\left(\frac{\beta -1}{2 - \beta}\right) + 1 \right\rceil \geq Ni\log_2(\beta)+\log_2\left(\frac{\beta -1}{2 - \beta}\right) + 1$, for all $i \in \N$, Lemma \ref{lem:s is in J beta} implies that $\ssf \in J_\beta^\omega$.
\end{proof}
We now proceed to formally prove that the dynamical system indeed generates a $\beta$-expansion.

\subsection{\texorpdfstring{The dynamical system produces a valid $\beta$-expansion}{}}

\label{sec:control of X and Z}

In this section, we show that for $\beta \in (1,2)$ and $s \in [0,1]$, if we define $N$ and $(\Sigma_i)_{i \in \N}$ according to (\ref{eq:final value parameters}), and $\ssf = s_0s_1\ldots $ according to (\ref{eq:s i to generate beta expansion}), then $E_\beta^N(\ssf)$ is a $\beta$-expansion of $s$.

\begin{theorem}\label{thm:convergence rational dynamical system to beta expansion}
	Let $\beta \in (1,2)$, $(N,(\Sigma_i)_{i \in \N})$ be defined according to $(\ref{eq:final value parameters})$ and set $\Sigma_0 = 0$. Let $s \in [0,1]$ and define $\ssf = s_0s_1\ldots$ according to (\ref{eq:s i to generate beta expansion}). Then,
	\begin{equation}
		E_\beta^N(\ssf) \in \Sigma_\beta(s).
	\end{equation}
\end{theorem}
\begin{proof}
	Let $\beta \in (1,2)$, $(N,(\Sigma_i)_{i \in \N})$ be defined according to $(\ref{eq:final value parameters})$ and set $\Sigma_0 = 0$. Let $s \in [0,1]$ and define $\ssf = s_0s_1\ldots$ according to (\ref{eq:s i to generate beta expansion}). By Theorem \ref{thm:convergence rational dynamical system to beta expansion}, $E_\beta^N(\ssf)$ is well-defined. To prove that $E_\beta^N(\ssf) \in \Sigma_\beta(s)$, by (\ref{eq:definition beta expansion}) we need to show that $\delta_\beta(E_\beta^N(\ssf)) = s$. We will successively show 
	\begin{equation}\label{eq:prefixes of d theta converge to d theta}
		\limi{n} \delta_\beta(E_\beta^N(s_0s_1\ldots s_{n-1})) = \delta_\beta(E_\beta^N(\ssf))
	\end{equation}
	and
	\begin{equation}\label{eq:prefixes of d theta converge to s}
		\limi{n} \delta_\beta(E_\beta^N(s_0s_1\ldots s_{n-1})) = s,
	\end{equation}
	which combination leads to $\delta_\beta(E_\beta^N(\ssf)) = s$. We start by proving (\ref{eq:prefixes of d theta converge to d theta}). Let $n \in \N$. Note that
	\begin{align}
		E_\beta^N(\ssf)_{1:Nn} &\overset{(\ref{eq:inferring chunks of greedy expansion from R i E beta})}{=} \left[\coprod_{i=0}^\infty d_\beta^N\left(R_i + s_i\right)\right]_{1:Nn}\\
		&= \coprod_{i=0}^{n-1} d_\beta^N\left(R_i + s_i\right)\\
		&\overset{(\ref{eq:inferring prefix of chunks of greedy expansion from R i E beta})}{=} E_\beta^N(s_0s_1 \ldots s_{n-1}) \label{eq:d theta prefix is a prefix of d theta}
	\end{align}
	It follows that
	\begin{align}
		\delta_\beta(E_\beta^N(\ssf)) &= \delta_\beta(E_\beta^N(\ssf)_{1:Nn} E_\beta^N(\ssf)_{Nn+1:\infty})\\
		&\osetlem{\ref{lem:factorisation of beta expansion through concatenation two}}{=} \delta_\beta(E_\beta^N(\ssf)_{1:Nn}) + \beta^{-Nn} \delta_\beta(E_\beta^N(\ssf)_{Nn+1:\infty})\\
		&\overset{(a)}{\in} \delta_\beta(E_\beta^N(\ssf)_{1:Nn}) + \beta^{-Nn} I_\beta\\
		&\overset{(\ref{eq:d theta prefix is a prefix of d theta})}{=} \delta_\beta(E_\beta^N(s_0s_1 \ldots s_{n-1}))  + \beta^{-Nn}I_\beta,
	\end{align}
	where (a) is by $\delta_\beta(\xf) \in I_\beta$, for all $\xf \in \{0,1\}^\omega$. Hence, 
	\begin{equation}
		0 \leq \limi{n} \left(\delta_\beta(E_\beta^N(\ssf)) - \delta_\beta(E_\beta^N(s_0s_1 \ldots s_{n-1}))\right) \leq \limi{n} \frac{\beta^{-Nn}}{\beta-1} = 0,
	\end{equation}
	yielding (\ref{eq:prefixes of d theta converge to d theta}). We move to the proof of (\ref{eq:prefixes of d theta converge to s}), which immediately follows from (\ref{eq:R i definition beta r}). Fix $n \in \N$. The dynamical system has been defined so that if (\ref{eq:s i to generate beta expansion}) holds, then 
	\begin{equation}
		R_n \overset{(\ref{eq:R i definition})}{=} \beta^{Nn}\left(\delta_2(\gf_2(s)_{\Sigma_i:\Sigma_{i+1}}) - \delta_\beta(E_\beta^N(s_0s_1 \ldots s_{n-1}))\right). 
	\end{equation}
	As (\ref{eq:final value parameters}) holds, Theorem \ref{thm:condition on the input sequences so that s is in dom E} implies that $\sf \in \dom E\beta^N$, i.e. $R_n + s_n \in I_\beta$. In particular, we have $0 \leq R_n \leq \frac{1}{\beta-1}$, so that
	\begin{equation}
		0 \leq \frac{R_i}{\beta^{Ni}} \overset{(\ref{eq:R i definition})}{=} \delta_2(\gf_2(s)_{\Sigma_i:\Sigma_{i+1}}) - \delta_\beta(E_\beta^N(s_0s_1 \ldots s_{n-1}))  \overset{(\ref{eq:R i definition})}{=} \frac{R_i}{\beta^{Ni}} \leq \frac{1}{(\beta-1) \beta^{Ni}} \underset{n \to \infty}{\to} 0.
	\end{equation}
	Therefore,
	\begin{equation}
		\limi{n} \delta_\beta(E_\beta^N(s_0s_1 \ldots s_{n-1})) = \limi{n} \delta_2(\gf_2(s)_{\Sigma_i:\Sigma_{i+1}}) = s.
	\end{equation}
	This ends the proof.
	% Let $n \in \N$. 
	% \allowdisplaybreaks
	% \begin{align}
	% 	\delta_\beta(E_\beta^N(s_0s_1 \ldots s_{n-1})) &\overset{(\ref{eq:inferring prefix of chunks of greedy expansion from R i E beta})}{=} \delta_\beta\left(\coprod_{i=0}^{n-1} d^N_\beta\left(R_i +  s_i\right)\right)\\
	% 	&\osetlem{\ref{lem:factorisation of beta expansion through concatenation}}{=}\sum_{i=0}^{n-1} \beta^{-Ni}\delta_\beta\left(d^N_\beta\left(R_i +  s_i\right)\right)\\
	% 	&\overset{(\ref{eq:definition multiple digit map})}{=}\sum_{i=0}^{n-1} \beta^{-Ni}\delta_\beta\left(\gf_\beta\left(R_i +  s_i\right)_{1:N}\right)\\
	% 	&\overset{(\ref{eq:n iterations of greedy map is same as reading the first n bits})}{=} \sum_{i=0}^{n-1} \beta^{-Ni} \left( R_i +  s_i - \frac{1}{\beta^N}G^N_\beta\left(R_i +  s_i\right) \right)\\
	% 	&\overset{(\ref{eq:greedy map chunks as dynamical system})}{=} \sum_{i=0}^{n-1} \beta^{-Ni} \left(R_i +s_i - \frac{1}{\beta^N}R_{i+1}\right)\\
	% 	&\overset{(\ref{eq:s i to generate beta expansion})}{=} \sum_{i=0}^{n-1} \beta^{-Ni} \left(R_i + \frac{\beta^{Ni}}{2^{\Sigma_i}}\delta_2(\gf_2(s)_{\Sigma_i+1:\Sigma_{i+1}}) - \frac{1}{\beta^N}R_{i+1}\right)\\
	% 	&= \sum_{i=0}^{n-1} \beta^{-Ni} R_i + \sum_{i=0}^{n-1} \frac{1}{2^{\Sigma_i}}\delta_2(\gf_2(s)_{\Sigma_i+1:\Sigma_{i+1}}) - \sum_{i=0}^{n-1} \beta^{-N(i+1)} R_{i+1}\\
	% 	&= R_0 - \beta^{Nn}R_n + \sum_{i=0}^{n-1} \frac{1}{2^{\Sigma_i}}\delta_2(\gf_2(s)_{\Sigma_i+1:\Sigma_{i+1}})\\
	% 	&\osetlem{\ref{lem:factorisation of beta expansion through concatenation}}{=} R_0 - \beta^{Nn}R_n + \delta_2(\gf_2(s)_{1:\Sigma_n}). \label{eq:intermediate equation delta beta x n}
	% \end{align}
	% From the definition of $E_\beta^N$, $R_0 = 0$. Moreover, as $(N, (\Sigma_i)_{i \in \N})$ satisfy (\ref{eq:final value parameters}), by (\ref{eq:R i is in 0 1}) one has $R_n \in [0,1]$. Therefore,
	% \begin{equation}
	% 	0 \geq \limi{n} \delta_\beta(E_\beta^N(s_0s_1 \ldots s_{n-1})) - \delta_2(\gf_2(s)_{1:\Sigma_n}) \geq \limi{n} -\beta^{-Nn} = 0.
	% \end{equation}
	% As $\limi{n} \Sigma_n = \infty$, we then have
	% \begin{equation}
	% 	\limi{n} \delta_\beta(E_\beta^N(s_0s_1 \ldots s_{n-1})) = \limi{n}\delta_2(\gf_2(s)_{1:\Sigma_n}) = s, 
	% \end{equation}
	% thereby finishing the proof.
\end{proof}

\subsection{\texorpdfstring{Algorithmic complexity of the $\beta$-expansion generated by the dynamical system}{}}


We have established that for $\beta \in (1,2)$, and $s \in [0,1]$, $E_\beta^N(\ssf)$ constitutes a valid $\beta$-expansion, assuming that (\ref{eq:final value parameters}) and (\ref{eq:s i to generate beta expansion}) hold. It remains to show that the algorithmic complexity of the prefixes of $E_\beta^N(\ssf)$ can be inferred from the algorithmic complexity of the prefixes of the greedy expansion $\gf_2(s)$. However, we can establish such a statement for $\beta \in \Q$ only. Indeed, the equations of the dynamical system that generate $E_\beta^N(\ssf)$ involve using the function $G_\beta$, which in full generality cannot be cast into an effective algorithm by virtue of requiring infinite precision. Instead, we would need another dynamical system which copes with approximations of $\beta$ only. We will develop this dynamical system in the next section. In this paragraph, we conclude the proof of (\ref{eq:inequality to be proven}), which establishes a relationship between the algorithmic complexity of the respective prefixes of $E_\beta^N(\ssf)$ and $\gf_2(s)$, respectively, for $\beta \in (1,2) \cap \Q$ only.

\begin{lemma}\label{lem:weak inequality binary beta expansion Kolmogorov rational}
	Let $\beta \in (1,2)\cap \Q$, $s \in [0,1]$ and $(N,(\Sigma_{i \in \N}))$ be defined according to (\ref{eq:final value parameters}). Then, there exists $\xf \in \Sigma_\beta(s)$, such that
	\begin{equation}
		K[\xf|N n] \leq K[\gf_2(s)|\Sigma_n] + c_\beta, \ \ \text{for all} \ n \in \N.
	\end{equation}
	where $c_\beta > 0$ depends only on $\beta$.
\end{lemma}
\begin{proof}
	Let $\beta \in (1,2)\cap \Q$, $s \in [0,1]$, $(N,(\Sigma_{i \in \N}))$ satisfying (\ref{eq:final value parameters}), and $\ssf = s_0s_1\ldots$ defined according to (\ref{eq:s i to generate beta expansion}). By Theorem \ref{thm:convergence rational dynamical system to beta expansion}, $E_\beta^N(s_0s_1\ldots s_{n-1})$ is an $Nn$-prefix of $E_\beta^N(\ssf) \in \Sigma_\beta(s)$, for $n \in \N$. Since $\beta \in \Q$, (\ref{eq:greedy map chunks as dynamical system}) can be run on a digital computer because it requires only finite precision, and hence $E_\beta^N$ can be cast into the effective algorithm $\varphi : \{0,1\}^\ast \to \{0,1\}^\ast$ summarized in Algorithm \ref{alg:conversion binary expansion to beta expansion}. Note that $\varphi(\gf_2(s)_{1:\Sigma_n}) = E_\beta^N(s_0s_1\ldots s_{n-1}) = E_\beta^N(\ssf)_{1:Nn}$, for $n \in \N$. Therefore, by Lemma \ref{lem:if two sequences are connected through an algorithm, they have same kolmogorov complexity} there exists $c_\varphi > 0$ such that
	\begin{equation}
		K[\xf|N n] \leq K[\gf_2(s)|\Sigma_n] + c_\varphi, \ \ \text{for all} \ n \in \N.
	\end{equation}
	$\varphi$ depends only on $E_\beta^N$, which itself depends only on $\beta$ as $N = \lceil \log_\beta(2) \rceil$. Hence, $c_\varphi$ depends only on $\beta$.
	\begin{algorithm}[H]
		\caption{$\beta$-expansion from a binary expansion}\label{alg:conversion binary expansion to beta expansion}
		\begin{algorithmic}[1]
			\Require $n \in \N$, $x \in \{0,1\}^\ast$
			\State Check that $|x| \geq \Sigma_n$, if not raise an error.
			\For{$i = 0, \ldots, n-1$}
			\State $s_i \gets \frac{\beta^{Ni}}{2^{\Sigma_i}}\delta_2(x_{\Sigma_i+1:\Sigma_{i+1}})$
			\EndFor
			\State \Return $E_\beta^N(s_0s_1\ldots s_{n-1})$
		\end{algorithmic}
	\end{algorithm}
\end{proof}

\begin{corollary}
	Let $\beta \in (1,2)\cap \Q$, $s \in [0,1]$. Then, there exists $\xf \in \Sigma_\beta(s)$, and $c_\beta > 0$ depending only on $\beta$, such that
	\begin{equation}\label{eq:algorithmic complexity of beta expansion lower than binary expansion}
		K[\xf|n] \leq K[\gf_2(s)|\lceil n \log_2(\beta)\rceil] + c_\beta, \ \text{for all} \ n \in \N.
	\end{equation}
\end{corollary}
\begin{proof}
	Let $\beta \in (1,2)\cap \Q$, $s \in [0,1]$, and $(N,(\Sigma_{i \in \N}))$ be defined according to (\ref{eq:final value parameters}). By Lemma (\ref{lem:weak inequality binary beta expansion Kolmogorov rational}), there exists $\xf \in \Sigma_\beta(s)$, such that 
	\begin{equation}
		K[\xf|N n] \leq K[\gf_2(s)|\Sigma_n] + C_\beta, \ \ \text{for all} \ n \in \N.
	\end{equation}
	where $C_\beta > 0$ depends only on $\beta$. By (\ref{eq:simple inequality on concatenation of sequences}), there exists $c_N>0$, depending only on $N$, such that
\begin{equation}\label{eq:before ultimate equation proof kolmogorov}
	K[\xf)|n]  \leq K[\gf_2(s)|\Sigma_{\lfloor n/N\rfloor}] + C_\beta + c_N, \ \ \text{for all} \ n \in \N.
\end{equation}
Moreover, by setting $M = 3 + \left| \log_2\left(\frac{\beta - 1}{2 - \beta}\right)\right|$, (\ref{eq:final value parameters}) yields
\begin{equation}
	\lceil n \log_2(\beta) \rceil - M \leq\Sigma_{\lfloor n/N\rfloor} \leq \lceil n \log_2(\beta) \rceil + M, \ \text{ for all } \ n \in \N.
\end{equation}
Then, by Lemma \ref{lem:kolmogorov complexity of closed sequences}, there exists $c_M > 0$, depending only on $M$, such that
\begin{equation}\label{eq:before ultimate equation proof kolmogorov 2}
	K[\gf_2(s)|\Sigma_{\lfloor n/N\rfloor}] \leq K[\gf_2(s)|\lceil n \log_2(\beta) \rceil] + c_M, \ \text{for all} \ n \in \N.
\end{equation}
By combining (\ref{eq:before ultimate equation proof kolmogorov}) and (\ref{eq:before ultimate equation proof kolmogorov 2}), one has
\begin{equation}
	K[\xf|n] \leq K[\gf_2(s)|\lceil n \log_2(\beta) \rceil] + C_\beta + c_N + c_M, \ \text{for all} \ n \in \N.
\end{equation}
As $N$ and $M$ depend only on $\beta$, $c_N$ and $c_M$ depend only on $\beta$ as well. We then finish the proof by letting $c_\beta := C_\beta + c_N + C_M$.
\end{proof}
We have established a simple relation between the algorithmic complexity of the prefixes of a certain $\beta$-expansion and the algorithmic complexity of the prefixes of the greedy binary expansion, for $\beta\in \Q$. The proof was effected through the study of a dynamical system that delivers the $\beta$-expansion. In the next section, we slightly modify the dynamical system $E_\beta^N$ to establish a similar relation, that holds for general $\beta \in (1,2)$.